\centerline{\bf \Huge Теория чисел. Добавка.} 

1. Пусть $n = {p_1}^{\alpha_1}\cdot{p_2}^{\alpha_2}\cdot{p_3}^{\alpha_3}\cdot\dotsc\cdot{p_k}^{\alpha_k}$, где $p_1$, $p_2$, ... , $p_k$ какие-то простые числа. Такое разложение называется каноническим разложением числа на степени простых. Докажите равенство:\\
$\varphi(n) = n(1 - \dfrac{1}{p_1})(1 - \dfrac{1}{p_2})\dots(1 - \dfrac{1}{p_k})$\\

Таким образом иногда легче считать функцию Эйлера от каких-то чисел, зная лишь какие простые числа входят в его разложение. Например как посчитать $\varphi(10000)$? Мы знаем, что в разложение 10000 входят много двоек и много пятерок, других простых делителей у него нет. Тогда $\varphi(10000) = 10000\cdot(1 - \dfrac{1}{2})(1 - \dfrac{1}{5}) = 10000\cdot \dfrac{1}{2} \cdot \dfrac{4}{5} = 5000 \cdot \dfrac{4}{5} = 4000$. \\

2. Для каких $n$ возможны равенства:\\
а) $\varphi(n) = n - 1$ \\ 
б) $\varphi(2n) = 2\varphi(n)$ \\ 
в) $\varphi(n^k) = n^{k-1}\varphi(n)$\\

3. Выпишем в ряд все правильные дроби со знаменателем n и сделаем возможные сокращения. Например, для $n = 12$ получится следующий ряд чисел: \\ \\
\centerline{$\dfrac{0}{1}, \ \dfrac{1}{12}, \ \dfrac{1}{6}, \
            \dfrac{1}{4}, \ \dfrac{1}{3}, \ \dfrac{5}{12}, \
            \dfrac{1}{2}, \ \dfrac{7}{12}, \ \dfrac{2}{3}, \
            \dfrac{3}{4}, \ \dfrac{5}{6}, \ \dfrac{11}{12}$} \\ 
            
Сколько получится дробей со знаменателем $d$, если $d$ — некоторый делитель числа $n$? \\ 

4. \textbf{Тождество Гаусса.} Докажите равенство: \\ 

\centerline{\larger$\sum\limits_{d|n}\varphi(d) = n$}
где знак $\sum\limits_{d|n}$ означает, что суммирование идет по всем делителям числа n. Например, посчитаем для $n = 12$. Делители 12 это 1, 2, 3, 4, 6 и 12. Значит, считаем $\varphi(1)+\varphi(2)+\varphi(3)+\varphi(4)+\varphi(6)+\varphi(12) = 1+1+2+2+2+4=12$
